\documentclass{article}
\usepackage{graphicx} % Required for inserting images

\usepackage[colorlinks=true, allcolors=blue]{hyperref}

\usepackage{amsmath}
\usepackage{amsfonts}
\usepackage{listings}
\usepackage{lstlangcoq}
\usepackage{matlab}
\lstset{mathescape=true,language=Coq}


\begin{document}
\title{Verifiable Circuit Wave Analysis}
\author{Timothy Welch and Andrew W. Appel}


\maketitle

\section{Introduction}

A \emph{rotating detonation engine} burns fuel (and oxidizer) in an
annular channel, that is, in the narrow space between two concentric
cylinders.  Detonation waves travel clockwise or counterclockwise
around the annulus, burning fuel faster (and therefore generating more
thrust) than if detonation was not occurring.

But are the detonation waves stable and long-lived?  With what speed
and direction do they travel?  Is there just one wave or several
simultaneously?  What is the shape of the wave?  How much combustion
that occurs is unrelated to the detonation waves?  To answer these
questions, Chacon, Feleo, and Gamba \cite{chacon21} aimed a camera at
the back of the engine, imaging the brightness of chemiluminescence at
101 equally spaced points around the cylinder, every 18 microseconds,
for durations of 1.5 seconds (thus, about 82,000 images).  Under the
hypothesis that in any 20ms timeframe there are just a few detonation
waves present, they used Fourier analysis to estimate the amplitude
and angular velocity of the detonation waves.  This analysis is called
\emph{circuit wave analysis} (CWA).  CWA takes into account that
the wave velocity may change (somewhat) over the timeframe, which allows
it to get significantly better results than a simple Fourier analysis
that would assume a constant velocity.
However, although Chacon \emph{et al.}'s CWA
deduces important properties of the circuit waves, it does not
guarantee that the waves it identifies account for all the activity in
the combustion annulus, nor does it \emph{describe} the gradual changes in velocity
of the waves, nor does it describe the shape of the waveform.

Based on Chacon \emph{et al.}'s CWA, we have designed and implemented
a \emph{Verified CWA} (VCWA), that produces a description $D$ of the
wave activity seen by the imaging device along with an accuracy
tolerance $\tau$.  It makes this formally verifiable\footnote{soon to
be formally verified} guarantee: VCWA's description of the waves
accounts for, within the accuracy tolerance, \emph{all} the
brightness variation seen by the imager.  In any one run of VCWA, the
accuracy tolerance may be small (indicating that it found a fairly
complete description of the activity) or large (indicating either that
the brightness variations could not be described by just a few waves
or that it could not identify those waves).

Either way, VCWA makes a verifiable claim: if you synthesize the
waveforms described by $D$ and subtract the original image sequence,
the 1-norm of the result will be $\le \tau$.  To verify this claim, run
a simple program that performs this synthesis and subtraction.

That is, we have an \emph{analyzer,} that uses Fourier analysis and
optimization algorithms to produce a \emph{hypothesis} $(D,\tau)$, and
a \emph{verifier}.  The analyzer is implemented in MATLAB.  The
verifier program is much simpler than the analyzer---since all it does
is synthesize and sum some waveforms---and we have implementations in
MATLAB and in C.

One can even apply \emph{formal verification} to VCWA.  Formal
verification, in this context, is a \emph{proof} in mathematical logic
that a program satisfies its mathematical specification.  Since any
such proof must be with respect to the minute details of the
programming language's formal semantics, there are far too many
details in such proofs to reliably perform or check such proofs by
hand.  Therefore, formal program verification usually uses software
\emph{proof checkers} to check that every step of such proofs is
valid, and \emph{proof assistants} and/or \emph{automated theorem
provers} to automate many steps of the proof.

We formally verify the verifier, not the analyzer.  We believe the analyzer is correct, but it need not
be \emph{proved} correct, since any hypothesis it proposes can be independently checked by the verifier.
Furthermore, this arrangement allows continued improvements of the analyzer's
(to describe a set of waves that more accurately account for the input data) without
altering the verifier or its proof.

\section{Specification}
\newcommand\Nfr{\ensuremath{{N_\mathrm{fr}}}}
\newcommand\Nb{\ensuremath{{N_\mathrm{b}}}}
\newcommand\NB{\ensuremath{{N_\mathrm{B}}}}
The \emph{input data} is a time series $X: \mathbb{R}^\Nfr \times \mathbb{R}^\Nb$
of $\Nfr$ \emph{frames} each containing $\Nb$ \emph{brightness} values.
Typically $\Nfr$ might be approximately $10^5$ and (in the experimental apparatus used) $\Nb=101$.  
The brightness
values are understood as equally spaced around the circumference of a circle, so that
points $i$ and \mbox{$(i+1) \mod \Nb$} are adjacent in (angular) space.
Between frame $k$ and frame $k+1$ there is
time interval $\delta$, typically 20 microseconds.
Thus, $X_{k,i}$
represents the brightness at time $k\delta$, angle $\theta(i) = \frac{2\pi}{\Nb}i$.  

\paragraph{Hypothesis, first attempt.}
One might hypothesize that brightness $X_{k,i}$ is the sum of $W$ waves,
where the $w$th wave $(0\le w <W)$ is described by a velocity $v_w$, amplitude $A_w$, phase $\theta_{w,0}$,
and shaped like a Gaussian distribution of width $\sigma_w$.  Therefore we synthesize
\begin{equation}\label{Snaive}
S^\mathrm{naive}_{k,i} = \sum_{w=0}^{W-1} A_w \mathrm{exp}(- \frac{(\theta_{w,0} + v_w k \delta -
  \frac{2\pi}{\Nb}i)^2}{2 \sigma_w^2}).
\end{equation}
However, this does not lead to an accurate estimation of $X$.  The waves do not have
constant amplitude, nor constant velocity, nor constant width.  Instead, all three of these
are slowly varying functions of time.\footnote{If these values were arbitrary chaotic functions of time,
it would be inappropriate to speak of a ``wave.''  The limitation to slowly varying functions is
intended to produce a hypothesis that has some scientific interpretability for the
engineer designing and measuring the combustion chamber.}

Furthermore, we found that there is a \emph{static background}.  That is, each of the $\Nb$ points around
the ring is consistently brighter or dimmer, regardless of $t$.  This may be a calibration issue, or
it may be that some parts of the ring are consistently hotter than others; either way, we can formulate
a hypothesis with a static background $B: \mathbb{R}^\Nb$, which is also a slowly varying function of time.

\paragraph{Hypothesis, final form.}
For $0 \le k < \Nfr$ and $0 \le i < \Nb$,
we explain $X_{k,i}$ as the sum of $W$ waves plus a static background.

To represent that the $w$th wave's amplitude, width, velocity are \emph{slowly varying} functions of time,
we represent each of those functions by a time series over $N_w$ equally spaced points, where $N_w < \Nfr$,
spanning the time $0$ to $(\Nfr-1)\delta$.
That is, for each $j$ such that $0 \le j \le N_w$,
the wave's amplitude at time  $\frac{\Nfr \delta}{N_w}k$ is
$A_{w,j}$, its width is $\sigma_{w,j}$,
and its velocity is $v_{w,j}$, and its initial position is $\theta_{w,0}$.

Then we use linear interpolation to reconstruct the value at time $k\delta$.
For example, the interpolation of $A$ at time $k \delta$
is calculated as follows.
Let $f=(\Nfr-1)/(N_w-1)$, the ratio between the real time range and the compressed index range.
Let $j=\lfloor k/f \rfloor$ and $i+1$ be the points
in (the compressed) $A$ between which we want to interpolate.
Then the result is
\begin{equation}\label{Ahat}
 \hat{A}_{w,k} = A_{w,j} + \frac{k-fj}{f}(A_{w,j+1}-A_{w,j}).
\end{equation}

The velocity wave $w$ is a slowly varying function of time $v_{w,j}$ where $j$ ranges
from $0$ to $N_w-1$.  We interpolate as above to calculate $\hat{v}$,
and we reconstruct the angular position of the $w$th wave at the $k$th time-step, as follows.
\begin{equation}\label{thetawk}
\theta_{w,k} = \left\{ 
\begin{array}{ll}
\theta_{w,0} & \mathrm{if}~k=0 \\ 
(\theta_{w,k-1}+\hat{v}_{w,k}\delta) \mod 2\pi  ~~~& \mathrm{if}~k>0
\end{array}
\right.
\end{equation}

The shape of wave $w$ is assumed to be a Gaussian with width $\sigma_{w,k}$ at time $k\delta$,
so the height of $w$ at time-step $k$ and position $i$ is,
\begin{equation}\label{awki}
a_{w,k,i} = \hat{A}_{w,i}\exp(-\mathrm{wrap}(2\pi\frac{i}{\Nb} - \theta_{w,k})^2 / (2\sigma ^2 ))
\end{equation}
where $\mathrm{wrap}(\phi)$ range-reduces $\phi$ to a value between $-\pi$ and $+\pi$.

The static background is represented as a two-dimensional array $B_{j,i}$,
for $0 \le j < \NB$ and $0 \le i < \Nb$, where $\NB < \Nfr$ to indicate that
$B$ is a slowly varying function of time.
However, in reconstruction we do not use linear interpolation, just simple rounding to nearest.
That is, the synthesized background at time $k$, position $i$ is,
\begin{equation}\label{Bki}
\hat{B}_{k,i} = B_{j,i}\qquad \mathrm{where}~j=\frac{\NB-1}{\Nfr-1}k
\end{equation}


Thus we define the synthesized brightness function as,
\begin{equation}\label{Ski}
S_{k,i} = \hat{B}_{k,i} + \sum_{w=0}^{W-1} a_{w,k,i}
\end{equation}

Now we can calculate the accuracy of this hypothesis.  The error $\epsilon$ is the
sum of the differences at each point:

\begin{equation}\label{epsilon}
\epsilon = \sum_{k=0}^{\Nfr-1}\sum_{i=0}^{\Nb-1} | S_{k,i}-X_{k,i} |
\end{equation}

\paragraph{Conclusion:}  To explain a dataset $(\Nfr,\Nb,X)$, a \emph{hypothesis} is a tuple $(D,\tau)$ where
$D=(B,(N_0,A_0,v_0,\sigma_0),\ldots (N_{w-1},A_{w-1},v_{w-1},\sigma_{w_1}))$;
where each $A_w$, $v_w$, and $\sigma_w$ is in $\mathbb{R}^{N_w}$;
and where the error $\epsilon = ||S(D)-X ||_1$ is $\le \tau $.

\section{How to find a hypothesis}

it ain't trivial.  Tim will explain \ldots

\section{How to verify a hypothesis}

Listing~\ref{Csynth} shows the C program of the verifier.

\begin{lstlisting}[language=C,float,label=Csynth,caption={The verifier as a C program}]
#define L (2.0*M_PI)
double range2pi(double x) {       // $\mathrm{calculate}~x \mod 2\pi $
  return x - L * floor (x / L);
}

double wrap(double x) {
  return range2pi(x+L/2.0)-L/2.0;
}

double interpolate(double f, double *y, int n, double k) {
    int j = floor(k/f);     $\mbox{\quad // equation (\ref{Ahat})}$
    return y[j] + ((k - j*f) / f) * (y[j+1] - y[j]);
}

void synthesize(Data *data, double *S) {
  int Nw = data->num_waves, Nb = data->Nb, Nfr = data->Nfr, Nbg  = data->Nbg;
  double dt = data->dt;
  Wave *waves = data->waves;
  double theta [ MAX_WAVES ];
  int w,k,i;

  for (k = 0; k < Nfr; k++) {
    for (i = 0; i < Nb; i++) {
      double a = data->BG[i + (int)(k*((double)(Nbg - 1)) / (Nfr - 1))*Nb];  $\mbox{// equation (\ref{Bki})}$

      for (w = 0; w < Nw; w++) {
        int Nt = waves[w].Nt;
        double f = (double)(Nfr - 1) / (double)(Nt - 1);
        double A     = interpolate(f, waves[w].A, Nt, k);
        double sigma = interpolate(f, waves[w].sigma, Nt, k);

	if (k==0)           $\mbox{// equation (\ref{thetawk})}$
	  theta[w] = waves[w].theta0;
	else if (i==0) {
	  double v = interpolate(f, waves[w].v, Nt, k-1);
	  theta[w] = range2pi (theta[w] + v * dt);
	}

        double dx = wrap(i * (L / Nb) - theta[w]);   
        a += A * exp(-dx*dx/(2.0*sigma*sigma));   $\mbox{// eqn. (\ref{awki})}$
      }
      S[i + k*Nb] = a;
} } }
\end{lstlisting}

\section{Proving correctness of the verifier}

not done yet 

\bibliographystyle{plain}
\bibliography{CWA}


\end{document}
